
\documentclass[12pt]{article}

\usepackage[utf8]{inputenc}
\usepackage{latexsym,amsfonts,amssymb,amsthm,amsmath,graphicx}

\setlength{\parindent}{0in}
\setlength{\oddsidemargin}{0in}
\setlength{\textwidth}{6.5in}
\setlength{\textheight}{8.8in}
\setlength{\topmargin}{0in}
\setlength{\headheight}{18pt}



\title{Math 075 Homework 5}
\author{Pranav Jayakumar}

\begin{document}
\maketitle
\vspace{1in}
\subsection*{Exercise 2.5.1}
\subsubsection*{Problem}
For each of the following elementary matrices,
describe the corresponding elementary row operations.\\
\subsubsection*{Part (a)}
\[
  E = \begin{bmatrix}
    1 & 0 & 3\\ 
    0 & 1 & 0\\ 
    0 & 0 & 1
  \end{bmatrix}
  \hspace{0.1in}
\]
\subsubsection*{Solution (a)}
\[
  E^{-1} = \begin{bmatrix}
    1 & 0 & \frac{1}{3}\\ 
    0 & 1 & 0\\ 
    0 & 0 & 1\\ 
  \end{bmatrix}
\]
\vspace{2in}
\subsection*{Exercise 2.5.2}
\subsubsection*{Problem}
In each case find an elementary matrix
$E$ such that $E = BA$. 
\subsubsection*{Part (a)}
\[
  A = \begin{bmatrix}
    2 & 1\\ 
    3 & -1
  \end{bmatrix},
  \hspace{0.1in}
  B = \begin{bmatrix}
    2 & 1\\ 
    1 & -2 
  \end{bmatrix}
\]
\subsubsection*{Solution (a)}
\[
  A \longrightarrow B
\]
\vspace{0.25in}
\[
  B = EA \longrightarrow \begin{bmatrix}
    2 & 1\\ 
    1 & -2 
  \end{bmatrix} = \begin{bmatrix}
    1 & 0\\ 
    -1 & 1 
  \end{bmatrix}\begin{bmatrix}
    2 & 1\\ 
    3 & -1
  \end{bmatrix}
\]
\[
  E = \begin{bmatrix}
    1 & 0\\ 
    -1 & 0 
  \end{bmatrix}
\]
\vspace{2in}
\subsection*{Exercise 2.5.3}
\subsubsection*{Problem}
Let $A = \begin{bmatrix}
  1 & 2\\ 
  -1 & 1 
\end{bmatrix}
$ and $C = \begin{bmatrix}
  -1 & 1\\ 
  2 & 1 
\end{bmatrix}
$. 
\subsubsection*{Part (a)}
Find elementary matrices $E_1$ and $E_2$ such that $C = E_2E_1A$
\subsubsection*{Solution (a)}
\[
  \begin{bmatrix}
    1 & 2\\ 
    -1 & 1 
  \end{bmatrix} \overset{R_2 \leftrightarrow R_1}{\longrightarrow}\begin{bmatrix}
    -1 & 1\\ 
    1 & 2 
    \end{bmatrix}\overset{R_2 \rightarrow R_2 - R_1}{\longrightarrow}\begin{bmatrix}
    -1 & 1\\ 
    2 & 1 
  \end{bmatrix}
\]
\[
  C = E_2E_1A \longrightarrow C = \begin{bmatrix}
    1 & 0\\ 
    1 & 1 
    \end{bmatrix}\begin{bmatrix}
    0 & 1\\ 
    1 & 0 
    \end{bmatrix}\begin{bmatrix}
    1 & 2\\ 
    -1 & 1 
  \end{bmatrix}
\]
\vspace{2in}
\subsection*{Exercise 2.5.8}
\subsubsection*{Problem}
In each case factor $A$ as a product of elementary matrices. 
\subsubsection*{Part (a)}
\[
  A = \begin{bmatrix}
    1 & 1\\ 
    2 & 1 
  \end{bmatrix}
\]
\subsubsection*{Solution (a)}
\[
  E_1 = \begin{bmatrix}
    0 & 1\\ 
    1 & 0 
  \end{bmatrix}\longrightarrow E_1(A) = \begin{bmatrix}
    2 & 1\\ 
    1 & 1 
  \end{bmatrix}
\]
\[
  E_2 = \begin{bmatrix}
    1 & -1\\ 
    0 & 1 
  \end{bmatrix}\longrightarrow E_2(E_1A) = \begin{bmatrix}
    1 & 0\\ 
    1 & 1 
  \end{bmatrix}
\]
\[
  E_3 = \begin{bmatrix}
    1 & 0\\ 
    -1 & 1 
  \end{bmatrix}\longrightarrow E_3(E_2E_1A) = \begin{bmatrix}
    1 & 0\\ 
    0 & 1 
  \end{bmatrix}
\]
\[
  E_3E_2E_1A = I 
\]
\[
  A = E^{-1}_1E^{-1}_2E^{-1}_3I
\]
\[
  A = \begin{bmatrix}
    0 & 1\\ 
    1 & 0 
  \end{bmatrix}\begin{bmatrix}
    1 & 1\\ 
    1 & 0 
  \end{bmatrix}\begin{bmatrix}
    1 & 0\\ 
    1 & 1 
  \end{bmatrix}\begin{bmatrix}
    1 & 0\\ 
    0 & 1 
  \end{bmatrix}
\]
\vspace{2in}
\subsection*{Exercise 2.6.1}
\subsubsection*{Problem}
Let $T: \mathbb{R}^3 \rightarrow \mathbb{R}^2$ be a linear transformation.
\subsubsection*{Part (a)}
Find $T\begin{bmatrix}
  8\\ 
  3\\ 
  7 
\end{bmatrix}
$ if: 
\[
  T\begin{bmatrix}
    1\\ 
    0\\ 
    -1 
  \end{bmatrix}
   = \begin{bmatrix}
     2\\ 
     3 
    \end{bmatrix}
\]
\[
  T\begin{bmatrix}
    2\\ 
    1\\ 
    3 
  \end{bmatrix} = \begin{bmatrix}
  -1\\ 
  0 
\end{bmatrix}
\]
\subsubsection*{Solution (a)}
Let:
\[
  \textbf{v}_1 = \begin{bmatrix}
    1\\ 
    0\\ 
    1 
  \end{bmatrix}, \hspace{0.1in} \textbf{v}_2 = \begin{bmatrix}
  2\\ 
  2\\ 
  3 
\end{bmatrix}
\]
\[
  \begin{bmatrix}
    8\\ 
    3\\ 
    7 
  \end{bmatrix} = a\textbf{v}_1 + b\textbf{v}_2 
\]
\[
  \begin{bmatrix}
    8\\ 
    3\\ 
    7 
  \end{bmatrix} = a\begin{bmatrix}
  1\\ 
  0\\ 
  1 
  \end{bmatrix} + b\begin{bmatrix}
  2\\ 
  1\\ 
  3 
  \end{bmatrix}
\]
\[
  \begin{cases}
    a + 2b = 8\\ 
    0a + b = 3\\ 
    a + 3b = 7 
  \end{cases}
\]
\[
  \begin{bmatrix}
    a\\ 
    b 
  \end{bmatrix} = RREF\left(\begin{bmatrix}
    1 & 2 & 8\\ 
    0 & 1 & 3\\ 
    -1 & 3 & 7 
  \end{bmatrix}\right)\longrightarrow\begin{bmatrix}
    1 & 0 & 2\\ 
    0 & 1 & 3\\ 
    0 & 0 & 0 
  \end{bmatrix}
\]
\[
  T\begin{bmatrix}
      8\\ 
      3\\ 
      7 
    \end{bmatrix} = T(2\textbf{v}_1 + 3\textbf{v}_2) = T(2\textbf{v}_1) + T(3\textbf{v}_2)
\]
\[
  T\begin{bmatrix}
    8\\ 
    3\\ 
    7 
  \end{bmatrix} = 2\begin{bmatrix}
  2\\ 
  3 
  \end{bmatrix} + 3\begin{bmatrix}
  -1\\ 
  0 
  \end{bmatrix}
\]
\[
  T\begin{bmatrix}
    8\\ 
    3\\ 
    7 
    \end{bmatrix} = \begin{bmatrix}
    4\\ 
    6 
    \end{bmatrix} + \begin{bmatrix}
    -3\\ 
    0 
  \end{bmatrix}
\]
\[
  T\begin{bmatrix}
    8\\ 
    3\\ 
    7 
  \end{bmatrix} = \begin{bmatrix}
  1\\ 
  6 
\end{bmatrix}
\]
\vspace{2in}
\subsection*{Exercise 2.6.2}
\subsubsection*{Problem}
Let $T: \mathbb{R}^4 \rightarrow \mathbb{R}^3$ be a linear transformation.
\subsubsection*{Part (a)}
Find $T\begin{bmatrix}
  1\\ 
  3\\ 
  -2\\ 
  -3 
\end{bmatrix}
$ if:
\[
  T\begin{bmatrix}
    1\\ 
    1\\ 
    0\\ 
    -1 
  \end{bmatrix} = \begin{bmatrix}
  2\\ 
  3\\ 
  1 
  \end{bmatrix}, \hspace{0.1in} T\begin{bmatrix}
  0\\ 
  -1\\ 
  1\\ 
  1 
\end{bmatrix}
\]
\subsubsection*{Solution (a)}
Let:
\[
  \textbf{v}_1 = \begin{bmatrix}
    1\\ 
    1\\ 
    0\\ 
    1 
  \end{bmatrix}, \hspace{0.1in} \begin{bmatrix}
  0\\ 
  -1\\ 
  1\\ 
  1 
\end{bmatrix}
\]
\[
  \begin{bmatrix}
    1\\ 
    3\\ 
    -2\\ 
    -3 
  \end{bmatrix} = a\textbf{v}_1 + b\textbf{v}_2 
\]
\[
  \begin{cases}
    a + 0b = 1\\ 
    a - b = 3\\ 
    0a + b = -2\\ 
    -a + b = -3 
  \end{cases}
\]
\[
  \begin{bmatrix}
    a\\ 
    b 
  \end{bmatrix} = RREF\left(\begin{bmatrix}
  1 & 0 & 1\\ 
  1 & -1 & 3\\ 
  0 & 1 & -2\\ 
  -1 & 1 & -3 
  \end{bmatrix}\right) \longrightarrow \begin{bmatrix}
  1 & 0 & 1\\ 
  0 & 1 & -2\\ 
  0 & 0 & 0\\ 
  0 & 0 & 0 
  \end{bmatrix} \longrightarrow \begin{bmatrix}
  1\\ 
  -2 
\end{bmatrix}
\]
\[
  T\begin{bmatrix}
    1\\ 
    3\\ 
    -2\\ 
    -3 
  \end{bmatrix} = T(\textbf{v}_1 - 2\textbf{v}_2) = T(\textbf{v}_1) + 2T(\textbf{v}_2)
\]
\[
  T\begin{bmatrix}
    1\\ 
    3\\ 
    -2\\ 
    -3 
  \end{bmatrix} = \begin{bmatrix}
  2\\ 
  3\\ 
  -1 
  \end{bmatrix} + \begin{bmatrix}
  10\\ 
  0\\ 
  2 
  \end{bmatrix}
\]
\[
  T\begin{bmatrix}
    1\\ 
    3\\ 
    -2\\ 
    -3 
  \end{bmatrix} = \begin{bmatrix}
  12\\ 
  3\\ 
  1 
  \end{bmatrix}
\]
\vspace{2in}
\subsection*{Exercise 2.6.3}
\subsubsection*{Problem}
In each case assume that the transformation T is linear, and use Theorem 2.6.2 to obtain the matrix A of T. 
\subsubsection*{Part (a)}
\[
  T: \mathbb{R}^2 \rightarrow \mathbb{R}^2 \hspace{0.1in} \text{is the reflection in the line} \hspace{0.1in} y = -x 
\]
\subsubsection*{Solution (a)}
\[
  T\begin{bmatrix}
    1\\ 
    0 
  \end{bmatrix} = \begin{bmatrix}
  0\\ 
  -1 
  \end{bmatrix}, \hspace{0.1in} T\begin{bmatrix}
  0\\ 
  1 
  \end{bmatrix} = \begin{bmatrix}
  -1\\ 
  0 
  \end{bmatrix}
\]
\[
  A = \begin{bmatrix}
    T\begin{bmatrix}
      1\\ 
      0 
    \end{bmatrix} & T\begin{bmatrix}
      1\\ 
      0 
    \end{bmatrix}
    \end{bmatrix} \longrightarrow A = \begin{bmatrix}
    0 & -1\\ 
    -1 & 0 
  \end{bmatrix}
\]
\vspace{2in}
\subsection*{Exercise 2.6.7}
\subsubsection*{Problem}
In each case show that $T: \mathbb{R}^2 \rightarrow \mathbb{R}^2$ is not a linear transformation. 
\subsubsection*{Part (a)}
\[
  T\begin{bmatrix}
    x\\ 
    y 
  \end{bmatrix} = \begin{bmatrix}
  xy\\ 
  0 
  \end{bmatrix}
\]
\subsubsection*{Solution (a)}
Let:
\[
  \textbf{u} = \begin{bmatrix}
    x_1\\ 
    y_1 
  \end{bmatrix}, \hspace{0.1in} \textbf{v} = \begin{bmatrix}
  x_2\\ 
  y_2 
  \end{bmatrix}
\]
$T$ is a linear transformation if:
\[
  T(\textbf{u} + \textbf{v}) = t(\textbf{u}) + T(\textbf{v})
\]
\[
  T\left(\begin{bmatrix}
    x_1\\ 
    y_1 
    \end{bmatrix} + \begin{bmatrix}
    x_2\\ 
    y_2 
  \end{bmatrix}\right) = T\begin{bmatrix}
  x_1\\ 
  y_1 
  \end{bmatrix} + T\begin{bmatrix}
  x_2\\ 
  y_2 
  \end{bmatrix}
\]
\[
  T\begin{bmatrix}
    x_1 + x_2\\ 
    y_1 + y_2 
  \end{bmatrix} = T\begin{bmatrix}
  x_1\\ 
  y_1 
  \end{bmatrix} + T\begin{bmatrix}
  x_2\\ 
  y_2 
  \end{bmatrix}
\]
\[
  \begin{bmatrix}
    (x_1 + x_2)(y_1 + y_2)\\ 
    0 
  \end{bmatrix} = \begin{bmatrix}
  x_1y_1\\ 
  0 
  \end{bmatrix} + \begin{bmatrix}
  x_2y_2\\ 
  0 
  \end{bmatrix}
\]
\[
  \begin{bmatrix}
    (x_1 + x_2)(y_2 + y_2)\\ 
    0 
  \end{bmatrix} \neq \begin{bmatrix}
  x_1y_1 + x_2y_2\\ 
  0 
  \end{bmatrix}
\]
$T$ does not satisfy the property of additivity. Therefore, it is not a linear transformation. 
\vspace{2in}
\end{document}
